% ============================================================
% BÖLÜM 7: HTP YORUMLAMA KURALLARI
% NIDA Bilgi Tabanının Temelini Oluşturan Akademik Çerçeve
% ============================================================

\chapter{HTP YORUMLAMA KURALLARI}

Bu bölüm, NIDA sisteminin bilgi tabanını oluşturan akademik yorumlama kurallarını sunmaktadır. Kurallar, temel kaynaklardan (Buck, 1948; Hammer, 1958; Jolles, 1971; Machover, 1949; Koch, 1952) derlenmiş ve güncel literatür ile desteklenmiştir.

\section{Ev (House) Çizimi Yorumlama Kuralları}

Buck'a (1948) göre ev çizimi, bireyin aile yaşamı, ev ortamı algısı ve güvenlik hissi hakkında bilgi vermektedir. Hammer (1958), evin ``psikolojik sığınak'' olarak yorumlanması gerektiğini vurgulamıştır.

\subsection{Çizim Boyutu ve Konumlandırma}

\begin{table}[H]
\centering
\caption{Ev Çizimi: Boyut ve Konumlandırma Yorumları}
\label{tab:evboyut}
\begin{tabular}{p{3.5cm}p{5cm}p{2.5cm}p{2.5cm}}
\toprule
\textbf{Gösterge} & \textbf{Yorum} & \textbf{Güven} & \textbf{Kaynak} \\
\midrule
Çok büyük ev & Aşırı telafi, ev ortamında baskı hissi & Orta & Hammer, 1958 \\
Çok küçük ev & Yetersizlik hissi, düşük benlik saygısı & Orta & Buck, 1948 \\
Üst konumlandırma & Hedeflere ulaşılamama, fantezi & Düşük & Jolles, 1971 \\
Alt konumlandırma & Güvensizlik, depresif eğilim & Düşük & Hammer, 1958 \\
Sol konumlandırma & Geçmişe yönelim, içe dönüklük & Düşük & Buck, 1948 \\
Sağ konumlandırma & Geleceğe yönelim, dışa dönüklük & Düşük & Buck, 1948 \\
\bottomrule
\end{tabular}
\end{table}

\subsection{Kapı Yorumlaması}

Kapı, sosyal etkileşim ve çevreye açıklık düzeyini temsil etmektedir (Hammer, 1958).

\begin{table}[H]
\centering
\caption{Ev Çizimi: Kapı Yorumları}
\label{tab:evkapi}
\begin{tabular}{p{3.5cm}p{5.5cm}p{2cm}p{2cm}}
\toprule
\textbf{Gösterge} & \textbf{Yorum} & \textbf{Güven} & \textbf{Kaynak} \\
\midrule
Kapı yok & Sosyal izolasyon, iletişim güçlüğü & Yüksek & Hammer, 1958 \\
Çok küçük kapı & Sosyal çekingenlik, isteksizlik & Orta & Buck, 1948 \\
Çok büyük kapı & Aşırı bağımlılık, onay arayışı & Orta & Jolles, 1971 \\
Açık kapı & Sıcaklık ve kabul ihtiyacı & Düşük & Hammer, 1958 \\
Kilitli/zincirli kapı & Savunmacılık, güvensizlik & Orta & Burns, 1987 \\
\bottomrule
\end{tabular}
\end{table}

\subsection{Pencere Yorumlaması}

Pencereler, çevre ile etkileşim ve algı açıklığını sembolize etmektedir (Buck, 1948).

\begin{table}[H]
\centering
\caption{Ev Çizimi: Pencere Yorumları}
\label{tab:evpencere}
\begin{tabular}{p{3.5cm}p{5.5cm}p{2cm}p{2cm}}
\toprule
\textbf{Gösterge} & \textbf{Yorum} & \textbf{Güven} & \textbf{Kaynak} \\
\midrule
Pencere yok & Çevreden kopukluk, içe kapanıklık & Yüksek & Buck, 1948 \\
Perdeli pencereler & Gizlilik ihtiyacı, savunmacılık & Orta & Hammer, 1958 \\
Parmaklıklı pencere & Hapsolmuşluk hissi, kısıtlanmışlık & Orta & Jolles, 1971 \\
Çok sayıda pencere & Çevreye açıklık, merak & Düşük & Buck, 1948 \\
Tek tarafta pencereler & Seçici algı, kısıtlı bakış açısı & Düşük & Hammer, 1958 \\
\bottomrule
\end{tabular}
\end{table}

\subsection{Çatı Yorumlaması}

Çatı, fantezi dünyası, entelektüel işlevler ve ego kontrolünü temsil etmektedir (Hammer, 1958).

\begin{table}[H]
\centering
\caption{Ev Çizimi: Çatı Yorumları}
\label{tab:evcati}
\begin{tabular}{p{3.5cm}p{5.5cm}p{2cm}p{2cm}}
\toprule
\textbf{Gösterge} & \textbf{Yorum} & \textbf{Güven} & \textbf{Kaynak} \\
\midrule
Orantısız büyük çatı & Fantezi dünyasına sığınma & Orta & Hammer, 1958 \\
Çatı yok & Düşünce yoksunluğu, somut düşünme & Orta & Buck, 1948 \\
Sivri çatı & Agresif eğilimler & Düşük & Jolles, 1971 \\
Düz çatı & Sınırlı hayal gücü & Düşük & Hammer, 1958 \\
Çatıda delik/hasar & Ego zayıflığı, kontrol kaybı & Orta & Burns, 1987 \\
\bottomrule
\end{tabular}
\end{table}

\subsection{Baca Yorumlaması}

Baca, duygusal sıcaklık ve aile dinamiklerini sembolize etmektedir (Hammer, 1958).

\begin{table}[H]
\centering
\caption{Ev Çizimi: Baca Yorumları}
\label{tab:evbaca}
\begin{tabular}{p{3.5cm}p{5.5cm}p{2cm}p{2cm}}
\toprule
\textbf{Gösterge} & \textbf{Yorum} & \textbf{Güven} & \textbf{Kaynak} \\
\midrule
Baca yok & Duygusal soğukluk, ilişki eksikliği & Düşük & Hammer, 1958 \\
Aşırı duman & Ev içi gerginlik, çatışma & Orta & Jolles, 1971 \\
Eğik baca & Dengesizlik, güvensizlik & Düşük & Buck, 1948 \\
\bottomrule
\end{tabular}
\end{table}

\section{Ağaç (Tree) Çizimi Yorumlama Kuralları}

Koch (1952) ve Bolander (1977), ağaç çiziminin bilinçdışı benlik algısını yansıttığını ileri sürmüştür. Ağaç, ``yaşam ağacı'' sembolizmi ile bireyin varoluşsal deneyimini temsil etmektedir.

\subsection{Gövde Yorumlaması}

Gövde, ego gücü ve temel benlik yapısını temsil etmektedir (Koch, 1952).

\begin{table}[H]
\centering
\caption{Ağaç Çizimi: Gövde Yorumları}
\label{tab:agacgovde}
\begin{tabular}{p{3.5cm}p{5.5cm}p{2cm}p{2cm}}
\toprule
\textbf{Gösterge} & \textbf{Yorum} & \textbf{Güven} & \textbf{Kaynak} \\
\midrule
İnce gövde & Yetersizlik, zayıflık hissi & Orta & Koch, 1952 \\
Kalın gövde & Güçlü ego, kararlılık & Orta & Bolander, 1977 \\
Yarık/delikli gövde & Travmatik deneyim & Yüksek & Hammer, 1958 \\
Eğik gövde & Dengesizlik, çevresel baskı & Orta & Koch, 1952 \\
Gövdede düğümler & Travma izleri, gelişimsel duraksama & Orta & Bolander, 1977 \\
\bottomrule
\end{tabular}
\end{table}

\subsection{Dal Yorumlaması}

Dallar, çevre ile ilişkiyi ve sosyal beceriler temsil etmektedir (Bolander, 1977).

\begin{table}[H]
\centering
\caption{Ağaç Çizimi: Dal Yorumları}
\label{tab:agacdal}
\begin{tabular}{p{3.5cm}p{5.5cm}p{2cm}p{2cm}}
\toprule
\textbf{Gösterge} & \textbf{Yorum} & \textbf{Güven} & \textbf{Kaynak} \\
\midrule
Yukarı uzanan dallar & İyimserlik, büyüme arzusu & Orta & Koch, 1952 \\
Aşağı sarkan dallar & Karamsarlık, depresif eğilim & Orta & Bolander, 1977 \\
Kırık dallar & Travma, kayıp deneyimi & Yüksek & Hammer, 1958 \\
Dal yok & Sosyal izolasyon, gelişim duraksaması & Yüksek & Koch, 1952 \\
Sivri/dikenli dallar & Agresif eğilimler, savunmacılık & Orta & Bolander, 1977 \\
\bottomrule
\end{tabular}
\end{table}

\subsection{Kök Yorumlaması}

Kökler, güvenlik hissi, temel ihtiyaçlar ve gerçeklik algısını temsil etmektedir (Koch, 1952).

\begin{table}[H]
\centering
\caption{Ağaç Çizimi: Kök Yorumları}
\label{tab:agackok}
\begin{tabular}{p{3.5cm}p{5.5cm}p{2cm}p{2cm}}
\toprule
\textbf{Gösterge} & \textbf{Yorum} & \textbf{Güven} & \textbf{Kaynak} \\
\midrule
Kök yok & Güvensizlik, köksüzlük hissi & Düşük & Koch, 1952 \\
Görünür kökler & Temel ihtiyaçların farkındalığı & Düşük & Bolander, 1977 \\
Aşırı büyük kökler & Güvenlik arayışı, bağımlılık & Orta & Hammer, 1958 \\
Ölü/çürük kökler & Umutsuzluk, varoluşsal kaygı & Orta & Koch, 1952 \\
\bottomrule
\end{tabular}
\end{table}

\section{İnsan (Person) Çizimi Yorumlama Kuralları}

Machover (1949), insan figürü çiziminin benlik algısı ve beden imajını yansıttığını ileri sürmüştür. Koppitz (1968), çocukların insan çizimlerinde duygusal göstergeler tanımlamıştır.

\subsection{Baş Yorumlaması}

Baş, entelektüel işlevleri, benlik algısını ve sosyal ilişkileri temsil etmektedir (Machover, 1949).

\begin{table}[H]
\centering
\caption{İnsan Çizimi: Baş Yorumları}
\label{tab:insanbas}
\begin{tabular}{p{3.5cm}p{5.5cm}p{2cm}p{2cm}}
\toprule
\textbf{Gösterge} & \textbf{Yorum} & \textbf{Güven} & \textbf{Kaynak} \\
\midrule
Büyük baş & Entelektüel odak, ego şişmesi & Orta & Machover, 1949 \\
Küçük baş & Yetersizlik, düşük özgüven & Orta & Koppitz, 1968 \\
Baş yok & Ciddi uyum sorunu (dikkat!) & Yüksek & Hammer, 1958 \\
Büyük gözler & Sosyal kaygı, şüphecilik & Orta & Machover, 1949 \\
Göz yok & Çevreden kaçınma, içe kapanıklık & Yüksek & Koppitz, 1968 \\
\bottomrule
\end{tabular}
\end{table}

\subsection{Gövde ve Uzuv Yorumlaması}

\begin{table}[H]
\centering
\caption{İnsan Çizimi: Gövde ve Uzuv Yorumları}
\label{tab:insangovde}
\begin{tabular}{p{3.5cm}p{5.5cm}p{2cm}p{2cm}}
\toprule
\textbf{Gösterge} & \textbf{Yorum} & \textbf{Güven} & \textbf{Kaynak} \\
\midrule
Kol yok & Yetersizlik, çaresizlik & Yüksek & Koppitz, 1968 \\
Bacak yok & Güvensizlik, hareket kısıtlılığı & Yüksek & Koppitz, 1968 \\
Uzun kollar & Çevreye ulaşma ihtiyacı & Düşük & Machover, 1949 \\
Kısa kollar & Çekingenlik, geri çekilme & Düşük & Machover, 1949 \\
El yok & Suçluluk, yetersizlik & Orta & Koppitz, 1968 \\
Büyük eller & Agresif eğilim, kompanzasyon & Orta & Machover, 1949 \\
Ayak yok & Güvensizlik, köksüzlük & Orta & Koppitz, 1968 \\
\bottomrule
\end{tabular}
\end{table}

\subsection{Koppitz Duygusal Göstergeleri}

Koppitz (1968), çocukların insan çizimlerinde 30 duygusal gösterge tanımlamıştır. Bu göstergelerden bazıları:

\begin{enumerate}
  \item Çarpık figürler
  \item Küçük figürler (sayfanın \%5'inden az)
  \item Figürlerin gizlenmesi
  \item Yoğun gölgeleme
  \item Cinsel organların belirginleştirilmesi
  \item Canavar/korkunç figürler
  \item Uzuv eksiklikleri
  \item Aşırı asimetri
\end{enumerate}

\section{Genel Yorumlama İlkeleri}

\subsection{Çizgi Kalitesi}

\begin{table}[H]
\centering
\caption{Genel Çizgi Kalitesi Yorumları}
\label{tab:cizgikal}
\begin{tabular}{p{3.5cm}p{5.5cm}p{2cm}p{2cm}}
\toprule
\textbf{Gösterge} & \textbf{Yorum} & \textbf{Güven} & \textbf{Kaynak} \\
\midrule
Hafif/zayıf çizgiler & Güvensizlik, çekingenlik & Orta & Hammer, 1958 \\
Ağır/koyu çizgiler & Gerginlik, agresyon & Orta & Hammer, 1958 \\
Kesik kesik çizgiler & Anksiyete, kararsızlık & Orta & Buck, 1948 \\
Titrek çizgiler & Anksiyete, motor kontrol sorunu & Orta & Machover, 1949 \\
\bottomrule
\end{tabular}
\end{table}

\subsection{Gölgeleme}

\begin{table}[H]
\centering
\caption{Gölgeleme Yorumları}
\label{tab:golge}
\begin{tabular}{p{3.5cm}p{5.5cm}p{2cm}p{2cm}}
\toprule
\textbf{Gösterge} & \textbf{Yorum} & \textbf{Güven} & \textbf{Kaynak} \\
\midrule
Yoğun gölgeleme & Anksiyete, obsesif eğilim & Orta & Koppitz, 1968 \\
Belirli bölgelerde & O bölgeye ilişkin kaygı & Orta & Machover, 1949 \\
Yüz gölgeleme & Benlik algısı sorunu & Yüksek & Koppitz, 1968 \\
\bottomrule
\end{tabular}
\end{table}

\section{Yaş ve Gelişim Evresine Göre Değerlendirme}

\begin{warningbox}
Yorumlamalar yapılırken çocuğun yaşı ve gelişim evresi mutlaka dikkate alınmalıdır. Lowenfeld'in (1947) gelişim evrelerine göre, bazı özellikler yaşa bağlı normal gelişimin parçası olabilir.
\end{warningbox}

\begin{table}[H]
\centering
\caption{Yaşa Göre Beklenen Çizim Özellikleri}
\label{tab:yas}
\begin{tabular}{cll}
\toprule
\textbf{Yaş} & \textbf{Beklenen Özellikler} & \textbf{Evre} \\
\midrule
5-6 & Kurbağa adam, belirsiz orantılar & Ön-Şematik \\
7-8 & Taban çizgisi, X-ray çizim & Şematik \\
9-10 & Detay artışı, örtüşme & Gerçekçilik \\
11-12 & Derinlik, perspektif denemeleri & Gerçekçilik \\
\bottomrule
\end{tabular}
\end{table}

\newpage
